% Created 2014-12-27 Sat 20:38
\documentclass[11pt]{article}
\usepackage[utf8]{inputenc}
\usepackage[T1]{fontenc}
\usepackage{fixltx2e}
\usepackage{graphicx}
\usepackage{longtable}
\usepackage{float}
\usepackage{wrapfig}
\usepackage[normalem]{ulem}
\usepackage{textcomp}
\usepackage{marvosym}
\usepackage{wasysym}
\usepackage{latexsym}
\usepackage{amssymb}
\usepackage{amstext}
\usepackage{hyperref}
\tolerance=1000
\author{Giulio Bottazzi}
\date{\today}
\title{gbget tutorial}
\hypersetup{
  pdfkeywords={},
  pdfsubject={},
  pdfcreator={Emacs 23.4.1 (Org mode 8.0.6)}}
\begin{document}

\maketitle
\tableofcontents


\section{Introduction to \texttt{gbget}}
\label{sec-1}

The special utility \texttt{gbget} is used to perform basic manipulation of
ASCII data files. It takes a file or a list of files containing
tabular data and extract command line specified parts of these
data. The result of the extraction can be further manipulated with a
list of possible transformations.

Essentially, the user provides a list of \emph{data specifications} at the
command line. The program analyze each specification in turns and,
according to it, extract and print data to standard output.

Each \emph{data specification} has the following structure

\begin{verbatim}
filename[block](col-range,rows-range)<options>
\end{verbatim}

where:

\begin{description}
\item[{filename}] is the name of a regular (ASCII) file

\item[{block}] identify a given data-block inside the file; data-blocks
are separated by 2 or more empty lines (format chosen for
consistency with gnuplot specifications). Notice that
comment lines (beginning with a "\#" symbol) DO NOT count as
empty lines for this purpose.

\item[{col-range}] the range of required columns, specified as
\verb~begin:end:skip~. Negative values are counted from the
end. Default is \verb~1:-1:1~ i.e. all columns. If begin>end
the columns are read in reversed order.
\end{description}


\begin{description}
\item[{rows-range}] the range of required rows, specified as
\verb~begin:end:skip~. Negative values are counted from the
end. Default is \verb~1:-1:1~ i.e. all rows. If begin>end
the rows are read in reversed order.

\item[{<options>}] is a list of single letter options that identify
successive transformations to be applied to data.
\end{description}

The list of options includes

\begin{description}
\item[{t}] transpose the matrix

\item[{f}] flatten the data column-wise, reducing them to a single column

\item[{l}] take the log of all fields

\item[{d}] take the column-wise difference: substract column 1 from 2,
column 2 from 3, etc.

\item[{D}] remove all lines containing at least one NAN entry

\item[{z}] remove the mean to each column

\item[{Z}] reduce each column to zscores, that is remove the mean and
divide by the standard deviation
\end{description}
\section{Examples}
\label{sec-2}

A few examples can help to understand the gbget syntax. Consider the
file \verb~test.dat~ with the following content

\begin{verbatim}
10 20 30
11 21 31
12 22 32
13 23 33
14 24 34


15 25 35
16 26 36
17 27 37
18 28 38
19 29 39
\end{verbatim}

i.e. two data blocks, separated by two blank lines, each made of three
columns and 5 rows. This file should already exists in the \verb~gbutils~
source directory. If the support for \href{http://www.gzip.org/zlib/}{zlib} has been found on your
system, in the examples below you can equivalently use the compressed
file \verb~test.dat.gz~.

Now if you type:

\begin{verbatim}
# gbget 'test.dat[1](1:2)'
\end{verbatim}

(the \textasciitilde{}'\textasciitilde{} are normally required to protect the content of the string
from shell expansion) you obtain

\begin{verbatim}
1.000000e+01  2.000000e+01
1.100000e+01  2.100000e+01
1.200000e+01  2.200000e+01
1.300000e+01  2.300000e+01
1.400000e+01  2.400000e+01
\end{verbatim}

i.e. the first two column \verb~(1:2)~ of the first datablock \verb~[1]~. By
default, the output is in scientific notation, but more on this below.

If instead you type

\begin{verbatim}
# gbget 'test.dat[2](-2:,2:3)'
\end{verbatim}

you obtain

\begin{verbatim}
2.600000e+01  3.600000e+01
2.700000e+01  3.700000e+01
\end{verbatim}

i.e. the rows from 2 to 3 (included) of the last two columns of the
second data block \verb~[2]~. Notice that a negative entry in column or row
specification means "count from the end".

You can also print \emph{each second column} of the second block with

\begin{verbatim}
# gbget 'test.dat[2](::2)'
\end{verbatim}

that gives you

\begin{verbatim}
1.500000e+01  3.500000e+01
1.600000e+01  3.600000e+01
1.700000e+01  3.700000e+01
1.800000e+01  3.800000e+01
1.900000e+01  3.900000e+01
\end{verbatim}

or \emph{each third row} of the whole file

\begin{verbatim}
# gbget 'test.dat(,::3)'
\end{verbatim}

to have

\begin{verbatim}
1.000000e+01  2.000000e+01  3.000000e+01
1.300000e+01  2.300000e+01  3.300000e+01
1.600000e+01  2.600000e+01  3.600000e+01
1.900000e+01  2.900000e+01  3.900000e+01
\end{verbatim}

If the initial and final positions in a slice specification are
reversed, the columns or the rows are printed in reverse order. For
instance

\begin{verbatim}
# gbget 'test.dat(,3:1)'
\end{verbatim}

gives

\begin{verbatim}
1.200000e+01  2.200000e+01  3.200000e+01
1.100000e+01  2.100000e+01  3.100000e+01
1.000000e+01  2.000000e+01  3.000000e+01
\end{verbatim}

Several different transformations can be applied to the chosen
\emph{matrix} of data.  For instance, the matrix can be flattened
column-wise, i.e. each column can be put after the previous one in a
single, long, column using the flag \verb~f~, or you can transpose it
using the option \verb~t~. Let see some examples. Consider

\begin{verbatim}
# gbget 'test.dat[1](,2:3)'
\end{verbatim}

which gives

\begin{verbatim}
1.100000e+01  2.100000e+01  3.100000e+01
1.200000e+01  2.200000e+01  3.200000e+01
\end{verbatim}

now you can flatten the output

\begin{verbatim}
# gbget 'test.dat[1](,2:3)f'
\end{verbatim}

\begin{verbatim}
1.100000e+01
1.200000e+01
2.100000e+01
2.200000e+01
3.100000e+01
3.200000e+01
\end{verbatim}

or transpose it

\begin{verbatim}
# gbget 'test.dat[1](,2:3)t'
\end{verbatim}

\begin{verbatim}
1.100000e+01  1.200000e+01
2.100000e+01  2.200000e+01
3.100000e+01  3.200000e+01
\end{verbatim}

Finally, the output format can be customized using options \verb~-o~ and
\verb~-e~. To fully exploit these options you need to know the syntax of
the \texttt{printf} command implemented in the standard C libraries. However,
the \emph{type} of output can be easily chosen with a single flag. If
instead of the scientific notation you prefer a fixed point notation
for your output, you have to use the option \verb~-o~ with the value \verb~%f~

\begin{verbatim}
# gbget 'test.dat[1](,2:3)t' -o ' %f'
\end{verbatim}

\begin{verbatim}
11.000000 12.000000
21.000000 22.000000
31.000000 32.000000
\end{verbatim}

while in order to have an \emph{integer} output you need \verb~-o~ with the
value \verb~%d~

\begin{verbatim}
# gbget 'test.dat[1](,2:3)t' -o ' %d'
\end{verbatim}

\begin{verbatim}
11 12
21 22
31 32
\end{verbatim}

in this case, however, be aware of the truncations: gbget rounds a
non-integer value down to the nearest integer!
% Emacs 23.4.1 (Org mode 8.0.6)
\end{document}